\documentclass[11pt]{article}

% Portable, dependency-free preamble — compiles with tectonic / pdflatex /
% latexmk / any LaTeX install, or directly on Overleaf. Matches the
% style of 01-mobe-bcn.tex.

\usepackage[letterpaper,margin=1in]{geometry}
\usepackage[utf8]{inputenc}
\usepackage[T1]{fontenc}
\usepackage{lmodern}
\usepackage{microtype}
\usepackage{hyperref}
\usepackage{url}
\usepackage{booktabs}
\usepackage{amsfonts}
\usepackage{amsmath}
\usepackage{amssymb}
\usepackage{mathtools}
\usepackage{graphicx}
\usepackage{xcolor}
\usepackage{natbib}
\usepackage{authblk}
\usepackage{titlesec}
\usepackage{abstract}
\usepackage{setspace}
\usepackage{listings}
\usepackage{fancyvrb}

% Typography
\hypersetup{
  colorlinks=true,
  linkcolor=blue!60!black,
  citecolor=blue!60!black,
  urlcolor=blue!50!black,
  pdftitle={Kazdov-LM: Scaling MoBE-BCN to Math Language Modeling via Cumsum Factorization},
  pdfauthor={Juan Cruz Dovzak},
}

\renewcommand{\abstractnamefont}{\normalfont\bfseries}
\renewcommand{\abstracttextfont}{\normalfont\small}

\titlespacing*{\section}{0pt}{1.4em}{0.6em}
\titlespacing*{\subsection}{0pt}{1.0em}{0.4em}

% Compact code listings (Python)
\lstdefinestyle{pycode}{
  basicstyle=\ttfamily\footnotesize,
  frame=single,
  framesep=4pt,
  framerule=0.3pt,
  rulecolor=\color{gray!40},
  backgroundcolor=\color{gray!5},
  breaklines=true,
  keepspaces=true,
  columns=fullflexible,
  aboveskip=0.5em,
  belowskip=0.5em,
  language=Python,
  commentstyle=\color{gray!70},
  keywordstyle=\color{blue!60!black},
  stringstyle=\color{red!50!black},
}
\lstset{style=pycode}

\colorlet{wipcolor}{orange!80!black}
\newcommand{\wip}[1]{{\color{wipcolor}\textbf{#1}}}
\newcommand{\TBD}{\wip{[TBD]}}

% -------------- Title --------------
\title{\textbf{Kazdov-LM: Scaling MoBE-BCN to Math Language Modeling\\
via Cumsum Factorization}}

\author{Juan Cruz Dovzak\thanks{Independent researcher, Buenos Aires, Argentina.
\texttt{juandovzak@gmail.com}. Code and data:
\url{https://github.com/OriginalKazdov/kazdov}}}

\date{April 2026 --- \wip{Work in Progress}}

\begin{document}

\maketitle

% -------------- WIP banner --------------
\begin{center}
\fbox{\parbox{0.93\textwidth}{\small\centering
\textbf{\color{wipcolor}Work-in-progress manuscript.} The companion paper
\citep{dovzak2026mobebcn} contains the complete MoBE-BCN architecture and
sample-efficiency results on algebraic composition tasks. Sections 1--3
of this manuscript (architecture, cumsum derivation, vectorized kernel)
are complete; Sections 4--7 (training results, ablations, qualitative
samples) will be finalised once the 5B-token alpha run completes.
Last updated: April 2026.
}}
\end{center}

% -------------- Abstract --------------
\begin{abstract}
\noindent
MoBE-BCN \citep{dovzak2026mobebcn} achieves state-of-the-art sample
efficiency on algebraic composition tasks but a naive implementation
materialises an $O(B \cdot L^2 \cdot D \cdot K)$ pairwise-composition tensor
inside attention, infeasible at language-model scale. We derive a
\textbf{cumsum-factorized attention kernel} that reduces MoBE-BCN attention
to $O(B \cdot L \cdot r \cdot K)$ when the router depends only on the
query side---aligning with standard token-level MoE routing and preserving
expert specialisation. Combined with a vectorized $K$-expert kernel, this
yields a $\mathbf{259\times}$ throughput improvement on a single
RTX~4090, enabling overnight training of a 98M-parameter math-specialised
LM (Kazdov-$\alpha$) on a 5.95B-token mathematics corpus.
On held-out math loss, Kazdov-$\alpha$ \TBD{} a matched-tokenizer
GPT-2~124M baseline. We release the model, code, and training recipe
under MIT.
\end{abstract}

% -------------- Introduction --------------
\section{Introduction}

Transformer attention is $O(L^2)$ in sequence length. Linear-attention
variants \citep{katharopoulos2020transformers,choromanski2020rethinking}
factor the softmax kernel into feature maps on $Q$ and $K$ that combine
via cumulative sums, yielding $O(L)$. MoBE-BCN's bilinear composition
$U((V^\top x) \odot (W^\top y))$ has precisely this structure---a tensor
product that factorises additively across the key axis---but the
naive implementation materialises the full pairwise tensor before
reducing. This paper shows the factorisation explicitly, derives the
$O(L)$ attention kernel, and demonstrates it enables math-specialised
language-model training on consumer hardware.

\paragraph{Contributions.}
\begin{enumerate}
\item \textbf{Cumsum-factorised MoBE-BCN attention} (Section~\ref{sec:cumsum}):
    a derivation and reference implementation that reduces memory from
    $O(B \cdot L^2 \cdot D \cdot K)$ to $O(B \cdot L \cdot r \cdot K)$ using
    query-side routing. Drop-in replacement for the naive causal kernel.
\item \textbf{Vectorized $K$-expert kernel} (Section~\ref{sec:vectorize}):
    stacked expert parameters with batched einsum; eliminates the
    Python per-expert loop.
\item \textbf{Kazdov-$\alpha$ training recipe} (Section~\ref{sec:setup}):
    98M-parameter math LM trained on 5.95B math tokens
    (OpenWebMath, proof-pile, MathPile, AutoMathText, NuminaMath) in
    $\sim$24h on a single RTX~4090.
\item \textbf{Held-out math-loss benchmark} (Section~\ref{sec:results},
    \wip{pending}): comparison against matched-tokenizer GPT-2 124M on
    five held-out math-corpus slices.
\item \textbf{Ablation study} (Section~\ref{sec:ablations},
    \wip{pending}): $K \in \{2, 4, 8\}$, $r \in \{64, 128, 256\}$,
    router type (Q-only vs Q+K), hybrid-MHA on/off.
\end{enumerate}

% -------------- Background --------------
\section{Background}

See the companion paper \citep{dovzak2026mobebcn} for the full MoBE-BCN
architecture. The core primitive is
\begin{align}
\text{BCN}_k(x, y) &= U_k \cdot \bigl((V_k^\top x) \odot (W_k^\top y)\bigr) + b_k,\\
\text{MoBE}(x, y) &= \sum_{k=1}^{K} w_k(x, y) \cdot \text{BCN}_k(x, y),
\end{align}
with $w_k(x, y) = \mathrm{softmax}(\text{Router}([x ; y]))_k$. In attention
this is applied as
\[
\text{out}_i \;=\; \frac{1}{n_{\mathrm{valid}}(i)} \sum_{j \leq i} \text{MoBE}(Q_i, K_j).
\]

\subsection{Memory and compute complexity (naive)}

The naive implementation expands $Q$ and $K$ to $(B, L, L, D)$ via
\texttt{unsqueeze}-and-\texttt{expand} before calling the composition:

\begin{lstlisting}
Qi = Q.unsqueeze(2).expand(B, L, L, D)    # O(B L^2 D)
Kj = K.unsqueeze(1).expand(B, L, L, D)    # O(B L^2 D)
fwd = MoBE(Qi, Kj)                        # O(B L^2 D K)
\end{lstlisting}

For a target 98M-parameter model ($D = 512$, $K = 4$, 12 layers) at $L = 256$,
each layer must hold $\sim\!4$~GB of intermediate activations. Combined
with gradient storage and the $O(L^2)$ softmax-based parallel MHA channel,
this exceeds 24~GB VRAM and is infeasible to train without gradient
checkpointing penalties that halve throughput.

% -------------- Cumsum factorisation --------------
\section{Cumsum-Factorised Attention}
\label{sec:cumsum}

\subsection{Key insight---bilinearity implies separability}

The $j$ index is the one we sum over in causal attention. Because
compose is \emph{bilinear} in $Q$ and $K$, the sum over $j$ factors
through the Hadamard product:
\begin{align}
\sum_{j \leq i} \text{BCN}(Q_i, K_j)
&= \sum_{j} U \cdot \bigl((V^\top Q_i) \odot (W^\top K_j)\bigr) \nonumber\\
&= U \cdot \bigl((V^\top Q_i) \odot \textstyle\sum_{j \leq i} (W^\top K_j)\bigr) \nonumber\\
&= U \cdot \bigl((V^\top Q_i) \odot \mathrm{cumsum}_j(W^\top K_j)_i\bigr).
\label{eq:factorization}
\end{align}
The $(V^\top Q_i)$ is constant in $j$ and factors out of the sum. The
$j$-cumulative sum of $W^\top K_j$ is a $(B, L, r)$ prefix-sum along the
sequence axis---a single pass, no pairwise materialisation. The
complexity drops from $O(B \cdot L^2 \cdot D \cdot K)$ to
$O(B \cdot L \cdot r \cdot K)$.

\subsection{Factorisability requires query-side routing}

Extending to MoBE, $\sum_j w_k(Q_i, K_j) \cdot \text{BCN}_k(Q_i, K_j)$.
If the routing weight $w_k(Q_i, K_j)$ depends on both $Q_i$ and $K_j$, it
does \emph{not} factor out of the sum over $j$, and the $O(L^2)$
materialisation returns. If the routing depends only on $Q_i$, it
factors:
\begin{align}
\sum_j w_k(Q_i) \cdot \text{BCN}_k(Q_i, K_j)
&= w_k(Q_i) \sum_j \text{BCN}_k(Q_i, K_j) \nonumber\\
&= w_k(Q_i) \cdot U_k \bigl((V_k^\top Q_i) \odot \mathrm{cumsum}_j(W_k^\top K_j)_i\bigr).
\end{align}
We therefore adopt \textbf{query-side routing}:
$w_k(Q_i) = \mathrm{softmax}(\text{Router}(Q_i))_k$. This aligns with
standard token-level MoE routing in modern
LLMs~\citep{shazeer2017outrageously,fedus2022switch,jiang2024mixtral}---the
router chooses experts per query token, not per $(Q, K)$ pair. Expert
specialisation survives in the per-token form; the algebraic results of
\citet{dovzak2026mobebcn} are compatible with this restriction (an
ablation on Power 2022 is reported in
Section~\ref{sec:ablations}).

\subsection{Vectorized $K$-expert kernel}
\label{sec:vectorize}

The naive MoBE implementation has a Python loop over experts:
\texttt{outs = torch.stack([e\_k(x, y) for e\_k in self.experts])}. This
serialises $K$ kernel launches and prevents GPU pipelining. We replace
it with stacked parameters and batched einsum:

\begin{lstlisting}
V_all = torch.stack([V_k for V_k in experts_V], dim=0)  # (K, D, r)
xV = torch.einsum('bld,kdr->blkr', x, V_all)            # (B, L, K, r)
yW = torch.einsum('bld,kdr->blkr', y, W_all)            # (B, L, K, r)
prod_kr  = (xV * yW).reshape(B, L, K * r)               # (B, L, K*r)
U_merged = U_all.permute(0, 2, 1).reshape(K * r, D)     # (K*r, D)
out      = prod_kr @ U_merged                           # (B, L, D)
\end{lstlisting}

A single batched matmul replaces $K$ sequential kernel launches.
Functionally identical (verified numerically, $\max |\Delta| < 10^{-7}$
vs naive CPU reference); throughput improves by $\sim\!3\times$ in
isolation.

\subsection{Combined: full kernel}

Putting the factorisation and vectorisation together, the causal
MoBE-BCN forward becomes:

\begin{lstlisting}
def forward(self, x):
    Q = self.W_Q(x); K = self.W_K(x)
    # Project per-expert  (B, L, K, r)
    xV = einsum('bld,kdr->blkr', Q, V_all)
    yW = einsum('bld,kdr->blkr', K, W_all)
    # Causal cumsum over j
    cum_yW = yW.cumsum(dim=1)                    # (B, L, K, r)
    prod   = xV * cum_yW                         # (B, L, K, r)
    # Q-side router
    w = softmax(self.router(Q), dim=-1)          # (B, L, K)
    prod_w = prod * w.unsqueeze(-1)              # (B, L, K, r)
    # Project to D via stacked U
    out = einsum('blkr,kdr->bld', prod_w, U_all) # (B, L, D)
    # Normalize by |valid j <= i|
    n = arange(1, L + 1, device=x.device).view(1, L, 1)
    return self.W_O(out / n)
\end{lstlisting}

Correctness verified against a naive pairwise+routing reference
implementation (max absolute error $< 10^{-7}$ on CPU; $< 10^{-5}$ under
bf16).

\subsection{Benchmark}
\label{sec:benchmark}

On a single RTX~4090 at model config $d = 512$, $L = 256$,
$n_{\mathrm{layers}} = 12$, $K = 4$, $r = 128$ (98M parameters total,
bf16 autocast, no gradient checkpointing):

\begin{table}[h]
\centering
\small
\begin{tabular}{lrrr}
\toprule
\textbf{Implementation} & \textbf{tok/s} & \textbf{VRAM peak} & \textbf{24h budget} \\
\midrule
Naive Python loop + pairwise expand & 288    & 15~GB & 25M \\
$+$ vectorized $K$ experts          & 853    &  8~GB & 74M \\
$+$ cumsum factorisation            & 6{,}006 & 1.3~GB & 519M \\
$+$ batch size 64 (memory freed)    & \textbf{74{,}654} & 19~GB & \textbf{6.5B} \\
\bottomrule
\end{tabular}
\caption{End-to-end throughput on a single RTX~4090. $\mathbf{259\times}$
total speedup. The $O(L)$ memory frees the batch-size budget, which
unlocks the final order-of-magnitude gain.}
\label{tab:benchmark}
\end{table}

% -------------- Training setup --------------
\section{Kazdov-$\alpha$ Training Setup}
\label{sec:setup}

\wip{Partial --- final numbers pending the complete training run.}

\paragraph{Model.}
KazdovLM 98M parameters --- $d = 512$, $L = 256$, $n_{\mathrm{layers}} = 12$,
$n_{\mathrm{heads}} = 8$, $r = 128$, $K = 4$, $\text{mlp\_dim} = 2048$,
$\text{max\_len} = 256$. Hybrid block: cumsum-factorised MoBE-BCN +
standard causal MHA in parallel.

\paragraph{Tokenizer.} GPT-2 BPE (50{,}257 vocab).

\paragraph{Data} (5.95B tokens):
\begin{itemize}
\item OpenWebMath 45\% (2.45B tokens)
\item proof-pile v1 22\% (1.40B)
\item MathPile-Commercial 15\% (1.05B)
\item AutoMathText top-quality slice 10\% (0.70B)
\item NuminaMath-CoT 8\% (0.35B)
\end{itemize}

\paragraph{Optimisation.}
AdamW $\text{lr} = 5 \times 10^{-4}$, $\beta = (0.9, 0.95)$,
$\text{wd} = 0.1$, grad clip $1.0$, WSD scheduler (warmup 500 steps,
stable, linear decay in last 20\%).

\paragraph{Hardware.}
Single RTX~4090, bf16 autocast, batch 32, length 256
$\rightarrow \sim\!60$--$70$K tok/s sustained.

\paragraph{Target.}
5B tokens $=$ $\sim\!610$K steps $\approx\!23$h wall time.

\paragraph{Held-out.}
Last 0.5\% of each data shard, deterministic eval.

% -------------- Results --------------
\section{Results}
\label{sec:results}

\wip{This section will be populated once training completes.}

\subsection{Held-out math LM loss}

\begin{table}[h]
\centering
\footnotesize
\setlength{\tabcolsep}{4pt}
\begin{tabular}{lrcrrrrrr}
\toprule
\textbf{Model} & \textbf{Params} & \textbf{Tokenizer} & \textbf{OWM} & \textbf{proof} & \textbf{MathPile} & \textbf{AutoMath} & \textbf{Numina} & \textbf{MEAN} \\
\midrule
GPT-2 124M (baseline) & 124M & GPT-2 & 3.57 & 3.12 & 3.04 & 3.28 & 2.58 & \textbf{3.12} \\
Kazdov-$\alpha$       & 98M  & GPT-2 & \TBD & \TBD & \TBD & \TBD & \TBD & \TBD \\
\bottomrule
\end{tabular}
\caption{\wip{Pending.} Held-out LM loss on 81{,}920 tokens per dataset
split (deterministic offsets). Lower is better. OWM: OpenWebMath;
proof: proof-pile v1.}
\label{tab:heldout}
\end{table}

\subsection{Loss curves}

\wip{Figure pending}---training-loss and held-out-loss curves vs.\
tokens seen.

\subsection{Routing diagnostics}

\wip{Figure pending}---per-layer expert usage entropy over training,
expert-specialisation heatmap.

% -------------- Ablations --------------
\section{Ablations}
\label{sec:ablations}

\wip{Pending.} Planned ablation sweep ($\sim\!200$M tokens each):

\begin{table}[h]
\centering
\small
\begin{tabular}{llll}
\toprule
\textbf{Ablation} & \textbf{Config} & \textbf{Held-out loss} & \textbf{$\Delta$ vs base} \\
\midrule
Base ($K = 4$, $r = 128$, hybrid-MHA) & --- & \TBD & 0 \\
$K = 2$                & half experts        & \TBD & \TBD \\
$K = 8$                & double experts      & \TBD & \TBD \\
$r = 64$               & narrower bilinear   & \TBD & \TBD \\
$r = 256$              & wider bilinear      & \TBD & \TBD \\
No hybrid MHA          & pure MoBE attention & \TBD & \TBD \\
Q+K routing (naive $O(L^2)$) & short-$L$ equivalence check & \TBD & \TBD \\
\bottomrule
\end{tabular}
\caption{Ablation plan. Rows will be populated from the D4 ablation sweep.}
\label{tab:ablations}
\end{table}

% -------------- Qualitative --------------
\section{Qualitative Samples}
\label{sec:samples}

\wip{Pending.} Prompt--generation pairs from the final checkpoint: LaTeX
completions, derivative/integral rules, short problem sketches.

% -------------- Discussion --------------
\section{Discussion}

\paragraph{Positioning.}
Kazdov-$\alpha$ is a demonstration of the cumsum-factorised MoBE-BCN
kernel at LM scale, not a production language model. At 98M parameters
and 5B tokens it is sub-Chinchilla-optimal by about a factor of 10 on
breadth; its purpose is architectural validation, not capability claims.
Downstream benchmarks (GSM8K, MATH) are not meaningfully evaluable at
this parameter count.

\paragraph{Relation to linear attention.}
The cumsum trick is the standard linearisation technique for any
attention whose kernel is separable across $Q$ and
$K$~\citep{katharopoulos2020transformers}. Our contribution is identifying
that (a) bilinear-composition attention is exactly such a kernel, and
(b) MoBE-BCN preserves separability precisely when routing is query-side,
which is the MoE convention anyway. The combination yields an $O(L)$
bilinear-composition attention with learned expert specialisation.

\paragraph{Limitations.}
\begin{itemize}
\item Short context (256 tokens in the alpha; scaling to 2048 is
    straightforward with the $O(L)$ kernel but not evaluated here).
\item Single-author, single-GPU training---reproducibility relies on the
    open release rather than independent confirmation.
\item No SFT or instruction-tuning; the model is base-LM only.
\item GSM8K/MATH benchmarks not meaningful at 98M; a 1B-scale follow-up
    is the correct target for capability claims.
\end{itemize}

% -------------- Acknowledgements --------------
\paragraph{Acknowledgments.}
This work was conducted independently without institutional support.

\bibliographystyle{plainnat}
\bibliography{references}

\end{document}
